% Vista científica exacta materializada desde una rama condicional.
% Fuente: source/demostraciones_primarias/31_06c_sectores_topologicos_no_omision.tex; rama: HMTIncludeTopologicalBlock.
% No modifica el contenido matemático de la rama de origen.
\chapter[Incidencia y sectores topológicos]{Incidencia, Fibonacci y sectores topológicos}
\label{chap:incidencia-fibonacci-topologia}
\index{sector topológico}\index{Fibonacci, anillo basado}
\index{grupo de trenzas}

La incidencia residuo--soporte prolonga el atlas hacia estructuras de fusión
y trenza. El orden es esencial: primero se construye el anillo basado
Fibonacci y sus dimensiones; después se estudian las representaciones
algebraicas de trenzas y el centro apuntado. Cada objeto conserva su propio
dominio y no se identifica con una partícula física por compartir nombre o
cardinal.

\section{Incidencia residuo--soporte}
\label{sec:incidencia-residuo-soporte}

Fijamos las clases de residuo y los soportes en el orden
\[
 \begin{aligned}
 C_1&=\{1,4,7\},&
 C_2&=\{2,5,8\},&
 C_0&=\{3,6,9\},\\
 \mathcal D_{\rm act}&=\{1,3,5,7\},&
 \mathcal D_{\rm evt}&=\{2,4,6,8\},&
 \mathcal D_{\bar0}^{(9)}&=\{9\}.
 \end{aligned}
\]
El último símbolo significa «soporte del residuo nonádico nulo». No designa
el dominio completo \(D_9=\{1,\ldots,9\}\) ni un sector físico de masa nula.

Sean \(E_{\mathcal D}\cong\mathbb Z^3\) y
\(E_C\cong\mathbb Z^3\) los módulos libres con bases ordenadas
\[
 (\mathcal D_{\rm act},\mathcal D_{\rm evt},
   \mathcal D_{\bar0}^{(9)})
 \quad\text{y}\quad
 (C_1,C_2,C_0),
\]
respectivamente. La aplicación de incidencia
\(B_{\rm res,sup}:E_{\mathcal D}\to E_C\) cuenta intersecciones.

\begin{theorem}[Forma de Smith e imagen de la incidencia]
\label{thm:incidencia-residuo-soporte}
En las bases precedentes,
\[
 B_{\rm res,sup}
 =\bigl(|C_i\cap\mathcal D_j|\bigr)_{i,j}
 =
 \begin{pmatrix}
 2&1&0\\
 1&2&0\\
 1&1&1
 \end{pmatrix}.
\]
Además,
\[
 \det B_{\rm res,sup}=3,\qquad
 \operatorname{SNF}(B_{\rm res,sup})
 =\operatorname{diag}(1,1,3),
\]
\[
 \operatorname{im}B_{\rm res,sup}
 =\{(u,v,w)\in\mathbb Z^3:u+v\equiv0\pmod3\}.
\]
Si
\[
 P=\begin{pmatrix}0&1&0\\1&0&0\\0&0&1\end{pmatrix},
\]
el intercambio simultáneo
\(C_1\leftrightarrow C_2\) y
\(\mathcal D_{\rm act}\leftrightarrow\mathcal D_{\rm evt}\) es equivariante:
\[
 P B_{\rm res,sup}P=B_{\rm res,sup}
 \qquad\Longleftrightarrow\qquad
 PB_{\rm res,sup}=B_{\rm res,sup}P.
\]
\end{theorem}

\begin{proof}
Las nueve entradas proceden de las intersecciones
\[
\begin{array}{cccc}
 &\mathcal D_{\rm act}&\mathcal D_{\rm evt}
 &\mathcal D_{\bar0}^{(9)}\\
 C_1&\{1,7\}&\{4\}&\varnothing\\
 C_2&\{5\}&\{2,8\}&\varnothing\\
 C_0&\{3\}&\{6\}&\{9\}.
\end{array}
\]
La expansión por la tercera columna da
\[
 \det B_{\rm res,sup}
 =\det\begin{pmatrix}2&1\\1&2\end{pmatrix}=3.
\]
El máximo común divisor de las entradas es uno. El menor de orden dos
formado por las filas primera y tercera y las columnas segunda y tercera
tiene determinante uno:
\[
 \det\begin{pmatrix}1&0\\1&1\end{pmatrix}=1.
\]
Por los divisores de Smith, \(d_1=d_2=1\). El producto
\(d_1d_2d_3\) coincide con el valor absoluto del determinante, que el
cálculo anterior fija en tres. Por consiguiente,

\[
 d_1d_2d_3=3.
\]
La forma normal es, por tanto, \(\operatorname{diag}(1,1,3)\).

Cada columna satisface \(u+v\equiv0\pmod3\), así que la imagen está contenida
en el submódulo indicado. Recíprocamente, si \(u+v\equiv0\pmod3\), entonces
\[
 x=\frac{2u-v}{3},\qquad
 y=\frac{2v-u}{3},\qquad
 z=w-x-y
\]
son enteros y
\[
 B_{\rm res,sup}(x,y,z)^{\mathsf T}=(u,v,w)^{\mathsf T}.
\]
Esto prueba la igualdad de imágenes. Finalmente, la multiplicación explícita
por \(P\) intercambia las dos primeras filas y las dos primeras columnas y
deja la matriz invariante.
\end{proof}

\begin{criterion}[Control de la incidencia]
\label{crit:incidencia-residuo-soporte}
Falsaría el teorema una sola intersección mal contada, un menor que alterase
los divisores de Smith, un vector de la imagen con \(u+v\not\equiv0\pmod3\),
un vector que satisficiera la congruencia pero no admitiese la preimagen
exhibida, o la pérdida de la equivarianza \(PBP=B\).
\end{criterion}

\section[Anillo basado de Fibonacci]{Anillo basado asociado a la incidencia
  \texorpdfstring{\(F_{\rm av}\)}{F\_av}}
\label{sec:anillo-basado-fibonacci}

El \cref{thm:descenso-necesario-fav} demuestra, a partir de la órbita de
modos \((\Sigma,\Pi,\Pi)\), la matriz
\[
 F_{\rm av}=
 \begin{pmatrix}0&1\\1&1\end{pmatrix}.
\]
De esta matriz se obtiene el siguiente anillo basado.

\begin{definition}[Anillo basado de Fibonacci]
\label{def:anillo-basado-fibonacci}
Sea
\[
 K_{\rm Fib}:=\mathbb Z\mathbf 1\oplus\mathbb Z\tau
\]
el grupo abeliano libre con unidad \(\mathbf 1\) y producto bilineal fijado por
\[
 \tau\cdot\mathbf 1=\tau,\qquad
 \tau\cdot\tau=\mathbf 1+\tau.
\]
En este estrato, \(+\) y \(\cdot\) son la adición y la multiplicación del
anillo. Los símbolos categóricos \(\oplus\) y \(\otimes\) se reservan para
una categorificación elegida posteriormente.
\end{definition}

\begin{theorem}[Realización de \(F_{\rm av}\) y dimensión de Perron]
\label{thm:anillo-fibonacci-fav}
En la base ordenada \((\mathbf 1,\tau)\), la multiplicación por \(\tau\) en
\(K_{\rm Fib}\) tiene matriz \(F_{\rm av}\). Si
\(F_0=0\), \(F_1=1\) y
\(F_{n+1}=F_n+F_{n-1}\) para \(n\geq1\), entonces
\[
 \tau^n=F_{n-1}\mathbf 1+F_n\tau
 \qquad(n\geq1).
\]
El único homomorfismo positivo de dimensión que envía
\(\mathbf 1\) a uno satisface
\[
 \operatorname{FPdim}(\tau)=\varphi.
\]
La identificación de las dos hojas tipadas con la base
\((\mathbf 1,\tau)\) determina por completo este anillo y su dimensión
positiva.
\end{theorem}

\begin{proof}
Las columnas de la matriz de multiplicación por \(\tau\) son las coordenadas
de
\[
 \tau\cdot\mathbf 1=\tau
 \quad\text{y}\quad
 \tau\cdot\tau=\mathbf 1+\tau;
\]
por tanto son \((0,1)^{\mathsf T}\) y \((1,1)^{\mathsf T}\), respectivamente.
Se obtiene así \(F_{\rm av}\).

La fórmula de potencias vale para \(n=1\). Si vale para \(n\), entonces
\[
 \tau^{n+1}
 =F_{n-1}\tau+F_n(\mathbf 1+\tau)
 =F_n\mathbf 1+F_{n+1}\tau,
\]
lo que cierra la inducción. Finalmente, una dimensión multiplicativa positiva
\(d=\operatorname{FPdim}(\tau)\) obedece
\[
 d^2=1+d.
\]
Sus raíces son \(\varphi\) y \(-\varphi^{-1}\); la positividad selecciona
\(d=\varphi\).
\end{proof}

La denominación \emph{Fibonacci} queda así explicada por la propia regla de
fusión: al iterar \(\tau\), los dos canales posibles aparecen con
multiplicidades consecutivas \(F_{n-1}\) y \(F_n\). El espacio de estados de
fusión crece, por tanto, con razón asintótica \(\varphi\). Ésta es la base
combinatoria de su uso computacional: la información se codifica en canales
globales de fusión y no en una masa asignada a un punto.

\begin{criterion}[Condiciones de una categorización Fibonacci]
\label{crit:categorificacion-fibonacci}
La capa anular quedaría refutada si la matriz de multiplicación no fuese
\(F_{\rm av}\), si fallase la recurrencia de potencias o si la dimensión
positiva no fuese \(\varphi\). Una categorización trenzada cuyo anillo de
Grothendieck sea el de la \cref{def:anillo-basado-fibonacci} exige además un
funtor monoidal
\[
 \mathfrak C_{\rm Fib}:\mathsf{TPK}
 \dashrightarrow\mathcal C_{\rm Fib},
\]
la verificación de pentágono y hexágono en su imagen y una regla interna que
seleccione calibre y quiralidad.
Estas condiciones no son consecuencias de la identidad anular por sí sola.
\end{criterion}

\section[Caracteres de trenza]{Seis caracteres y cuatro torsiones de trenza}
\label{sec:rescate-trenzas-ex02}

Sea \(B_n\), \(n\geq2\), el grupo de trenzas de Artin con generadores
\(\sigma_1,\ldots,\sigma_{n-1}\), relaciones
\[
 \sigma_i\sigma_{i+1}\sigma_i
 =\sigma_{i+1}\sigma_i\sigma_{i+1},
 \qquad
 \sigma_i\sigma_j=\sigma_j\sigma_i\quad(|i-j|\geq2),
\]
y sea \(w:B_n\to\mathbb Z\) la suma de exponentes. Los lectores
\[
 w_2:=w\bmod2,\qquad w_3:=w\bmod3
\]
se reúnen, por el teorema chino del resto, en
\[
 w_6:B_n\longrightarrow
 \mathbb Z/2\mathbb Z\times\mathbb Z/3\mathbb Z
 \simeq\mathbb Z/6\mathbb Z.
\]

\begin{theorem}[Abelianización, seis caracteres y torsiones de permutación]
\label{thm:trenzas-caracteres-z6}
Para \(n\geq2\),
\[
 B_n^{\rm ab}\simeq\mathbb Z,
\]
mediante \(w\). En consecuencia, \(w_6\) produce seis caracteres unitarios
\[
 \chi_k(b)
 :=\exp\!\left(\frac{2\pi\mathrm i k}{6}w(b)\right),
 \qquad k\in\mathbb Z/6\mathbb Z.
\]

Sea \(V\ne0\) un espacio de Hilbert, sea
\(P_i:V^{\otimes n}\to V^{\otimes n}\) el intercambio de los factores
\(i,i+1\), y sea \(q\in\mu_6\). La regla
\[
 \rho_q(\sigma_i):=qP_i,
 \qquad
 \rho_q(b)=q^{w(b)}\pi_n(b),
\]
donde \(\pi_n\) es la representación de permutación, define una
representación de \(B_n\). Esta representación desciende al grupo simétrico
si y sólo si
\[
 q^2=1.
\]
Para los cuatro valores \(q\in\mu_6\setminus\{\pm1\}\),
\[
 \operatorname{Inv}(\rho_q)=0.
\]
Esos cuatro valores definen torsiones escalares genuinamente trenzadas, pero
no constituyen por sí solos cuatro espacios de Fock ni cuatro sectores
anyónicos físicos.
\end{theorem}

\begin{proof}
En todo cociente abeliano, la relación de trenza da
\[
 2[\sigma_i]+[\sigma_{i+1}]
 =[\sigma_i]+2[\sigma_{i+1}],
\]
luego todos los generadores tienen la misma clase. No queda otra relación
sobre esa clase, pues \(w(\sigma_i)=1\); por tanto la abelianización es
\(\mathbb Z\).

Los operadores \(P_i\) satisfacen las relaciones de trenza y
\(P_i^2=I\). Multiplicar cada generador por el mismo escalar \(q\) conserva
las dos relaciones, de modo que \(\rho_q\) es una representación. Para
factorizar a través de \(S_n\) debe satisfacer la relación adicional
\(\sigma_i^2=1\), y
\[
 \rho_q(\sigma_i)^2=q^2I.
\]
Esto prueba el criterio \(q^2=1\).

Si \(v\) fuera invariante, entonces
\(qP_i v=v\), es decir, \(P_i v=q^{-1}v\), para todo \(i\). Como cada
involución \(P_i\) sólo tiene autovalores \(\pm1\), no existe tal vector no
nulo cuando \(q\ne\pm1\). Por consiguiente,
\(\operatorname{Inv}(\rho_q)=0\) para los cuatro valores restantes.
\end{proof}

Cuando \(\dim V>1\) y \(n\geq3\), los operadores \(P_1\) y \(P_2\) no
conmutan; por ello la representación completa no se vuelve abeliana por
proceder de un carácter escalar. El carácter \(\chi_k\) y la representación
\(\rho_q\) son objetos distintos: el primero es unidimensional; la segunda
retiene la acción de permutación.

\section[Centro modular y realizaciones topológicas]{Centro modular apuntado y condiciones de realización Fibonacci e Ising}
\label{sec:centro-apuntado-no-go}

Sea
\[
 A=(\mathbb Z/9\mathbb Z)^2
\]
y sea \(\operatorname{Vec}_A\) la categoría de espacios vectoriales
\(A\)-graduados con asociador trivial. Fijar un bicharácter no degenerado
identifica \(\widehat A\) con \(A\), pero no altera los conteos siguientes.

\begin{theorem}[Centro apuntado y obstrucción dimensional]
\label{thm:centro-apuntado-no-go}
El centro de Drinfeld
\[
 \mathcal Z(\operatorname{Vec}_A)
\]
es una categoría modular apuntada cuyos simples se parametrizan por
\((a,\chi)\in A\times\widehat A\). Por tanto posee
\[
 |A|\,|\widehat A|=81^2=6561
\]
simples, todos de dimensión cuántica uno, y dimensión cuántica total
\[
 \mathcal D
 =\sqrt{\sum_{x\ {\rm simple}}d_x^2}
 =\sqrt{6561}=81.
\]
En esta categoría no puede realizarse, preservando simples y dimensiones de
Frobenius--Perron, ni el anillo Fibonacci ni el sector Ising: éstos requieren
respectivamente un simple de dimensión \(\varphi\) y un simple de dimensión
\(\sqrt2\).
\end{theorem}

\begin{proof}
Para un grupo abeliano finito y asociador trivial, cada objeto simple del
centro queda determinado por un grado \(a\in A\) y un carácter
\(\chi\in\widehat A\) que especifica su media trenza. Tanto el objeto
graduado como el carácter son unidimensionales; en consecuencia, todos los
simples son invertibles y tienen dimensión cuántica uno. Como
\(|A|=|\widehat A|=9^2=81\), el número de pares es \(81^2\), y la fórmula de
\(\mathcal D\) da \(81\).

Toda subcategoría de fusión de una categoría apuntada vuelve a ser apuntada;
en particular, sus simples tienen dimensión uno. Esto contradice
\(\operatorname{FPdim}(\tau)=\varphi\) para Fibonacci y
\(d_\sigma=\sqrt2\) para Ising. La obstrucción es dimensional y no depende
de una semejanza de nombres o cardinales.
\end{proof}

\begin{criterion}[Condiciones para una realización física de los sectores]
\label{crit:promocion-sectores-topologicos}
Los resultados de
\cref{thm:trenzas-caracteres-z6,thm:centro-apuntado-no-go} quedarían
algebraicamente falsados por un fallo de las relaciones de Artin, del
criterio \(q^2=1\), del censo \(6561\), de la apuntabilidad o de
\(\mathcal D=81\). Una realización física exige, además, un
Hamiltoniano o protocolo impulsado, una brecha espectral, operadores de
trenza realizables, protección frente a perturbaciones y un mapa desde las
multisecciones enriquecidas. Ninguno de esos datos se deduce únicamente de
los dos teoremas algebraicos anteriores.
\end{criterion}
